\documentclass[11pt]{article}

\usepackage[margin=1in]{geometry}
\usepackage{amsmath,amssymb,amsthm}
\usepackage{enumitem}
\usepackage{hyperref}
\usepackage{xcolor}
\usepackage{booktabs}

\hypersetup{colorlinks=true, linkcolor=blue, urlcolor=blue, citecolor=blue}

\title{TOAE209/TOAH209 -- Design and Analysis of Algorithms\\[4pt]
\large Week 2: Theoretical Questions}
\author{Foreign Trade University}
\date{}

\begin{document}
\maketitle

\section*{Instructions}

Part A contains 20 questions requiring written explanations or short proofs. Part B is a
10-question quiz (true/false, multiple choice, and short answer). Attempt every question
on your own before consulting Part C (Solutions) at the end of this document.

\section*{Part A: Theory Questions}

\begin{enumerate}[label=\textbf{A\arabic*.}, leftmargin=2.2em]

\item State the formal ($c$, $n_0$) definitions of $f(n) = O(g(n))$, $f(n) = \Omega(g(n))$,
and $f(n) = \Theta(g(n))$. Explain in one sentence each what they mean informally (upper
bound, lower bound, tight bound).

\item Prove, directly from the definition, that $f(n) = 5n^2 + 3n + 2$ satisfies
$f(n) = O(n^2)$. Exhibit explicit constants $c$ and $n_0$.

\item Prove, directly from the definition, that $f(n) = 5n^2 + 3n + 2$ satisfies
$f(n) = \Omega(n^2)$. Exhibit explicit constants $c$ and $n_0$. Conclude that $f(n) =
\Theta(n^2)$.

\item Prove that Big-O is \emph{reflexive} ($f(n) = O(f(n))$) and \emph{transitive} (if
$f(n) = O(g(n))$ and $g(n) = O(h(n))$, then $f(n) = O(h(n))$).

\item Prove that $f(n) = \Omega(g(n))$ if and only if $g(n) = O(f(n))$.

\item True or false, with justification: if $f(n) = O(g(n))$, then $2^{f(n)} = O(2^{g(n)})$.
(Hint: consider $f(n) = 2n$ and $g(n) = n$.)

\item Order the following eight functions from slowest- to fastest-growing, and justify
the relative order of any two consecutive functions in your list:
\[
  n^3, \quad \log_2 n, \quad n!, \quad n\sqrt{n}, \quad 2^n, \quad n^2 \log n, \quad n,
  \quad 100.
\]

\item Is $n^2 = O(n^3)$? Is $n^3 = O(n^2)$? Is $n^2 = \Theta(n^3)$? For each, answer
yes/no and justify briefly using the definitions.

\item For a polynomial $p(n) = a_d n^d + \cdots + a_1 n + a_0$ with $a_d > 0$, prove that
$p(n) = O(n^d)$. (You do not need to prove the matching $\Omega(n^d)$ bound, though you may
for extra rigor.)

\item Consider the code fragment
\begin{verbatim}
total = 0
for i in range(n):
    for j in range(i):
        total += 1
\end{verbatim}
Derive an exact closed-form expression for the final value of \texttt{total} as a function
of $n$, and state its tightest $\Theta(\cdot)$ bound.

\item Consider the recurrence $T(n) = T(n-1) + 1$ for $n \ge 1$, with $T(0) = 1$. Solve it
in closed form (by unrolling) and state its $\Theta(\cdot)$ class. Which algorithmic
pattern does this recurrence correspond to?

\item Consider the recurrence $T(n) = T(\lfloor n/2 \rfloor) + 1$ for $n \ge 1$, with
$T(0)=1$. Solve it in closed form (assuming $n = 2^k$ for simplicity) and state its
$\Theta(\cdot)$ class. Which algorithmic pattern does this recurrence correspond to?

\item Consider the recurrence $T(n) = 2T(n/2) + n$ for $n > 1$, with $T(1) = 1$. Using the
recursion-tree method (sum the work done at each level of recursion), show that $T(n) =
\Theta(n \log n)$.

\item Define what it means for an algorithm to run in \emph{polynomial time}. Explain why
``polynomial time'' is considered a reasonable (if imperfect) threshold for ``efficient'',
giving one argument based on \emph{composability} of polynomial-time subroutines.

\item Give a concrete numerical example (a table of running times for $n=10, 30, 50$, say)
that illustrates the practical gap between an $O(n^3)$ algorithm and an $O(2^n)$
algorithm. At what rough value of $n$ does the $O(2^n)$ algorithm become impractical (say,
$> 10^9$ operations)?

\item State the loop invariant used to prove the correctness of iterative binary search
(\textsc{BinarySearchIterative}). Explain the three parts of a loop-invariant proof
(initialization, maintenance, termination) as applied to this invariant.

\item Binary search shrinks the candidate range by (at least) half on every iteration.
Use this fact to derive the $\Theta(\log n)$ bound on the number of iterations, being
careful about the ``at least half'' vs.\ ``exactly half'' distinction (the range size is
$\lceil m/2 \rceil$ or $\lfloor m/2 \rfloor$ depending on which half is kept, where $m$ is
the previous range size).

\item Explain the ``divide into pairs'' trick for finding both the max and the min of $n$
elements using roughly $\frac{3n}{2}$ comparisons instead of $2(n-1)$. Why are both
approaches $\Theta(n)$, and in what sense is the paired approach still an improvement?

\item Explain the amortized-analysis argument for why appending to a dynamic array that
doubles its capacity on overflow costs $O(1)$ \emph{amortized} per append, even though
individual appends that trigger a resize cost $\Theta(n)$. Your answer should mention the
total number of element copies over $n$ appends.

\item \textbf{(Tie-in to Week 1: Gale--Shapley complexity).} Week 1's dictionary-based
Gale--Shapley implementation used \texttt{.index()} to find a person's rank in a
preference list, where \texttt{.index()} on a Python list of length $n$ is $O(n)$. The
algorithm's main loop runs $O(n^2)$ times (Week 1, Lemma on proposals), and each iteration
may call \texttt{.index()} a constant number of times. What is the resulting worst-case
running time of this implementation, as a function of $n$? Now suppose (as in this week's
Problem 16) we precompute a rank table \texttt{women\_rank[w][m]} once, in $O(n^2)$ time,
so that each rank lookup becomes $O(1)$. What is the running time of the algorithm using
this precomputed table? Explain why precomputing the rank table is worth the $O(n^2)$
upfront cost.

\end{enumerate}

\newpage
\section*{Part B: Quiz}

\begin{enumerate}[label=\textbf{B\arabic*.}, leftmargin=2.2em]

\item \textbf{(True/False)} If $f(n) = \Theta(g(n))$, then $f(n) = O(g(n))$ and $f(n) =
\Omega(g(n))$.

\item \textbf{(Short answer)} What is the tightest $\Theta(\cdot)$ bound for $f(n) = 7
\log n + n^{0.5}$?

\item \textbf{(Multiple choice)} Which of the following is the \emph{tightest} valid
asymptotic bound for $f(n) = n^2 + 1000n$?
\begin{enumerate}[label=(\alph*)]
    \item $O(n)$
    \item $O(n \log n)$
    \item $\Theta(n^2)$
    \item $\Theta(n^3)$
\end{enumerate}

\item \textbf{(True/False)} $2^{n+1} = \Theta(2^n)$.

\item \textbf{(Short answer)} For the recurrence $T(n) = T(n/2) + 1$, $T(1) = 1$, what is
$T(1024)$?

\item \textbf{(Multiple choice)} A loop \texttt{for i in range(n): for j in range(i, n):
...} (with $O(1)$ work per inner iteration) runs the inner body a total of:
\begin{enumerate}[label=(\alph*)]
    \item $\Theta(n)$ times
    \item $\Theta(n \log n)$ times
    \item $\Theta(n^2)$ times
    \item $\Theta(n^3)$ times
\end{enumerate}

\item \textbf{(True/False)} An algorithm with worst-case running time $O(n^{50})$ is, by
the definition used in this course, a polynomial-time algorithm.

\item \textbf{(Short answer)} What is $\lceil \log_2(16+1) \rceil$, the worst-case number
of recursive calls for binary search on a sorted array of length $n=16$?

\item \textbf{(True/False)} For the dynamic-array doubling strategy, the total number of
element copies after $n = 17$ appends (starting from capacity 1) is strictly less than
$2n = 34$.

\item \textbf{(Short answer)} In the efficient array/queue-based Gale--Shapley
implementation (Problem 16), what is the time complexity of \texttt{build\_women\_rank},
as a function of $n$?

\end{enumerate}

\newpage
\section*{Part C: Solutions}

\subsection*{Part A}

\begin{enumerate}[label=\textbf{A\arabic*.}, leftmargin=2.2em]

\item $f(n) = O(g(n))$: there exist constants $c > 0, n_0 \ge 0$ such that $f(n) \le
c\,g(n)$ for all $n \ge n_0$ (\emph{$f$ grows no faster than $g$, an upper bound}).
$f(n) = \Omega(g(n))$: there exist $c > 0, n_0 \ge 0$ such that $f(n) \ge c\,g(n)$ for all
$n \ge n_0$ (\emph{$f$ grows at least as fast as $g$, a lower bound}). $f(n) =
\Theta(g(n))$: both of the above hold (\emph{$f$ and $g$ grow at the same rate, a tight
bound}).

\item Take $c = 6$ and $n_0 = 4$. For $n \ge 4$: $3n + 2 \le 3n + 2n = 5n \le n \cdot n =
n^2$ (using $3n+2 \le 5n$ for $n \ge 1$, and $5n \le n^2$ for $n \ge 5$; to be safe take
$n_0=5$). Then $f(n) = 5n^2 + (3n+2) \le 5n^2 + n^2 = 6n^2 = c\,n^2$ for $n \ge n_0 = 5$.
Hence $f(n) = O(n^2)$ with $c=6, n_0=5$.

\item Take $c = 5$ and $n_0 = 0$. For all $n \ge 0$, $f(n) = 5n^2 + 3n + 2 \ge 5n^2 = c
\cdot n^2$ (since $3n+2 \ge 0$). Hence $f(n) = \Omega(n^2)$ with $c=5, n_0=0$. Combined with
A2, $f(n) = O(n^2)$ and $f(n) = \Omega(n^2)$, so $f(n) = \Theta(n^2)$.

\item \textbf{Reflexivity:} take $c=1, n_0=0$; then $f(n) \le 1 \cdot f(n)$ trivially for
all $n \ge 0$, so $f(n) = O(f(n))$.

\textbf{Transitivity:} Suppose $f(n) = O(g(n))$, so there exist $c_1, n_1$ with $f(n) \le
c_1 g(n)$ for $n \ge n_1$. Suppose $g(n) = O(h(n))$, so there exist $c_2, n_2$ with $g(n)
\le c_2 h(n)$ for $n \ge n_2$. Let $n_0 = \max(n_1,n_2)$ and $c = c_1 c_2$. For $n \ge n_0$:
$f(n) \le c_1 g(n) \le c_1 (c_2 h(n)) = c\,h(n)$. Hence $f(n) = O(h(n))$.

\item ($\Rightarrow$) If $f(n) = \Omega(g(n))$, there exist $c>0, n_0$ with $f(n) \ge
c\,g(n)$ for $n \ge n_0$, i.e.\ $g(n) \le \frac{1}{c} f(n)$ for $n \ge n_0$. Taking $c' =
1/c > 0$, this is exactly the statement $g(n) = O(f(n))$. ($\Leftarrow$) Symmetric: if
$g(n) = O(f(n))$, there exist $c>0, n_0$ with $g(n) \le c\,f(n)$ for $n \ge n_0$, i.e.\
$f(n) \ge \frac{1}{c} g(n)$ for $n \ge n_0$, giving $f(n) = \Omega(g(n))$ with constant
$1/c$.

\item \textbf{False.} Take $f(n) = 2n$ and $g(n) = n$; clearly $f(n) = O(g(n))$ (with
$c=2$). But $2^{f(n)} = 2^{2n} = (2^n)^2 = (2^{g(n)})^2$. For $2^{2n} = O(2^n)$ we would
need a constant $c$ with $2^{2n} \le c \cdot 2^n$ for all large $n$, i.e.\ $2^n \le c$ --
impossible since $2^n \to \infty$. So $2^{f(n)} \ne O(2^{g(n)})$ in general: Big-O is
\emph{not} preserved under exponentiation.

\item From slowest- to fastest-growing:
\[
  100 \;\prec\; \log_2 n \;\prec\; n \;\prec\; n\sqrt{n} \;\prec\; n^2 \log n \;\prec\; n^3
  \;\prec\; 2^n \;\prec\; n!.
\]
Justifications: $100 \prec \log_2 n$ since $\log_2 n \to \infty$ while $100$ is constant.
$\log_2 n \prec n$ since $\log n = o(n)$. $n \prec n\sqrt{n} = n^{1.5}$ trivially (larger
exponent). $n^{1.5} \prec n^2 \log n$ since $n^2 \log n / n^{1.5} = n^{0.5}\log n \to
\infty$. $n^2 \log n \prec n^3$ since $n^3 / (n^2 \log n) = n/\log n \to \infty$. $n^3
\prec 2^n$ since every polynomial is $o(2^n)$. $2^n \prec n!$ since $n!/2^n =
\prod_{i=1}^n (i/2) \to \infty$ (eventually every factor exceeds $1$).

\item \textbf{$n^2 = O(n^3)$: yes.} Take $c=1, n_0=1$; for $n \ge 1$, $n^2 \le n^3$.
\textbf{$n^3 = O(n^2)$: no.} If it held, there would be $c, n_0$ with $n^3 \le c\,n^2$,
i.e.\ $n \le c$ for all $n \ge n_0$ -- impossible since $n \to \infty$. \textbf{$n^2 =
\Theta(n^3)$: no}, since $\Theta$ requires both $O$ and $\Omega$, and $n^2 = \Omega(n^3)$
would require $n^2 \ge c\,n^3$, i.e.\ $1/n \ge c$ for all large $n$ -- impossible for
$c>0$.

\item For $n \ge 1$, each term satisfies $|a_i| n^i \le |a_i| n^d$ (since $i \le d$ and $n
\ge 1$ implies $n^i \le n^d$). Hence $p(n) \le \sum_{i=0}^d |a_i| n^i \le \left(
\sum_{i=0}^d |a_i| \right) n^d$. Taking $c = \sum_{i=0}^d |a_i|$ and $n_0 = 1$ gives $p(n)
\le c\,n^d$ for all $n \ge 1$, so $p(n) = O(n^d)$.

\item The inner loop runs $i$ times for each value of $i = 0, 1, \ldots, n-1$, so
$\texttt{total} = \sum_{i=0}^{n-1} i = \frac{(n-1)n}{2} = \binom{n}{2}$. This is
$\Theta(n^2)$.

\item Unrolling: $T(n) = T(n-1)+1 = T(n-2)+2 = \cdots = T(0) + n = 1 + n$. So $T(n) = n+1 =
\Theta(n)$. This is the pattern of \textbf{linear recursion} -- e.g.\ a simple recursive
countdown or a single recursive pass over a list/array (one recursive call per element).

\item Assuming $n = 2^k$: $T(2^k) = T(2^{k-1})+1 = \cdots = T(2^0) + k = T(1) + k = 1+k =
1 + \log_2 n$. So $T(n) = \Theta(\log n)$. This is the pattern of \textbf{halving
recursion} -- e.g.\ binary search, or any algorithm that discards half its input at each
step.

\item Assume $n=2^k$. At recursion depth $i$ ($i=0,\ldots,k-1$), there are $2^i$
sub-problems, each of size $n/2^i$. The non-recursive (``$+n$'' at the top level, or
proportionally $+(n/2^i)$ per sub-problem at depth $i$) work at depth $i$ totals $2^i
\cdot \frac{n}{2^i} = n$. Summing over $k = \log_2 n$ levels (depth $0$ through $k-1$),
plus $\Theta(n)$ base cases at depth $k$ (each of size $1$, $T(1)=1$, and there are $n$ of
them), the total is $n \cdot \log_2 n + \Theta(n) = \Theta(n \log n)$.

\item An algorithm runs in \textbf{polynomial time} if its worst-case running time is
$O(n^k)$ for some constant $k$ independent of the input. \textbf{Composability argument:}
if algorithm $A$ runs in $O(n^a)$ and calls, as a subroutine, algorithm $B$ running in
$O(n^b)$, a polynomial (e.g.\ $O(n^c)$) number of times, the overall running time is
$O(n^c \cdot n^b) + O(n^a) = O(n^{b+c}) + O(n^a)$, still polynomial. So polynomial-time
algorithms can be freely composed without ``escaping'' the polynomial-time class --
unlike, say, composing two exponential algorithms (which could give a doubly-exponential
result).

\item
\begin{center}
\begin{tabular}{@{}lrrr@{}}
\toprule
Running time & $n=10$ & $n=30$ & $n=50$ \\
\midrule
$n^3$ & $1{,}000$ & $27{,}000$ & $125{,}000$ \\
$2^n$ & $1{,}024$ & $1{,}073{,}741{,}824$ & $\approx 1.13 \times 10^{15}$ \\
\bottomrule
\end{tabular}
\end{center}
At $n=30$, $O(n^3)$ is already over $10^4$ times smaller than $O(2^n)$. $O(2^n)$ exceeds
$10^9$ operations once $n \gtrsim 30$ (since $2^{30} \approx 1.07 \times 10^9$), while
$O(n^3)$ does not reach $10^9$ until $n \gtrsim 1000$. So the $O(2^n)$ algorithm becomes
impractical around $n \approx 30$--$40$, while the $O(n^3)$ algorithm remains usable for
$n$ in the thousands.

\item The invariant is: \emph{``if $x$ is present in $A$, its index lies in
$[lo, hi]$.''} \textbf{Initialization:} before the first iteration, $[lo,hi]=[0,n-1]$,
which trivially contains every index, so the invariant holds. \textbf{Maintenance:} if
$A[mid] \ne x$, sortedness of $A$ guarantees that if $x$ is present it must be on the side
of $mid$ consistent with the comparison $A[mid]$ vs.\ $x$; updating $lo$ or $hi$
accordingly keeps $x$'s index (if it exists) inside the new $[lo,hi]$. \textbf{Termination:}
the loop ends either by finding $x$ at $mid$ (correct), or when $lo > hi$ -- by the
invariant, an empty range means $x$ is not present, so returning $-1$ is correct.

\item Let $m_k$ be the range size before the $k$-th iteration, $m_0 = n$. Each iteration
checks $A[mid]$ and keeps one of the two halves, each of size at most $\lceil (m_k-1)/2
\rceil \le \lceil m_k/2 \rceil$, so $m_{k+1} \le \lceil m_k/2 \rceil$. Repeatedly halving
(rounding up) from $n$ reaches $0$ after at most $\lceil \log_2(n+1) \rceil$ steps -- more
precisely, $m_k \le \lceil n/2^k \rceil$, and the loop must terminate once $m_k = 0$, which
happens for $k \ge \log_2(n+1)$. Hence the iteration count is $\Theta(\log n)$, with the
precise worst case $\lceil \log_2(n+1) \rceil$.

\item For each pair $(a,b)$ of the $n$ elements (assume $n$ even for simplicity, $n/2$
pairs): 1 comparison decides which of $a,b$ is the pair's local max and which is the local
min. Then 2 more comparisons update the running max (against the local max) and running
min (against the local min). Total: $3 \cdot (n/2) = \frac{3n}{2}$ comparisons, versus the
naive $2(n-1) \approx 2n$ for two separate scans. Both are $\Theta(n)$ -- the asymptotic
class is unchanged -- but $\frac{3n}{2} < 2n$ for all $n>0$, so the paired approach uses
fewer comparisons by a constant factor (a real, if modest, improvement within the same
$\Theta$ class).

\item Resizes (doublings) happen when the size, just before an append, equals the current
capacity, i.e.\ at sizes $1, 2, 4, 8, \ldots, 2^{k-1}$ where $2^k \ge n$ is the final
capacity. The $i$-th resize (from capacity $2^{i-1}$ to $2^i$) copies $2^{i-1}$ elements.
Total copies $= 1+2+4+\cdots+2^{k-1} = 2^k - 1 < 2n$ (since $2^{k-1} < n \le 2^k$ implies
$2^k < 2n$). So the total cost of copying across all $n$ appends is $O(n)$, even though a
single resize costs $\Theta(n)$. Dividing the total $O(n)$ cost over $n$ appends gives
$O(1)$ \emph{amortized} cost per append -- the expensive resizes are ``paid for'' by the
many cheap appends between them.

\item With \texttt{.index()} lookups (each $O(n)$), and $O(n^2)$ total loop iterations
each doing $O(1)$ such lookups, the worst-case running time is $O(n^2) \times O(n) =
O(n^3)$. With a precomputed rank table (built in $O(n^2)$ time -- filling an $n \times n$
table, one entry per (woman, rank) pair), each lookup becomes $O(1)$, so the main loop
costs $O(n^2) \times O(1) = O(n^2)$, plus the $O(n^2)$ precomputation, for a total of
$O(n^2)$ -- asymptotically better than $O(n^3)$ for large $n$. The upfront $O(n^2)$ cost is
worth it because it is \emph{dominated by} (same order as, in fact smaller in practice
than) the main loop's cost, and it eliminates a per-iteration $O(n)$ factor that would
otherwise dominate everything.

\end{enumerate}

\subsection*{Part B}

\begin{enumerate}[label=\textbf{B\arabic*.}, leftmargin=2.2em]

\item \textbf{True.} This is exactly the definition of $\Theta$: $f(n)=\Theta(g(n))$ means
$f(n)=O(g(n))$ and $f(n)=\Omega(g(n))$.

\item $\Theta(n^{0.5})$ (equivalently $\Theta(\sqrt{n})$). The $\sqrt{n}$ term dominates
$\log n$ asymptotically (any positive power of $n$ dominates $\log n$), so the sum is
$\Theta(\sqrt{n})$.

\item \textbf{(c) $\Theta(n^2)$.} The $n^2$ term dominates the lower-order $1000n$ term;
$1000$ is just a constant absorbed into the $\Theta$.

\item \textbf{True.} $2^{n+1} = 2 \cdot 2^n$, so taking $c_1=c_2=2$ (and $n_0=0$) gives $2
\cdot 2^n \le 2^{n+1} \le 2 \cdot 2^n$ trivially (it's exactly $2\cdot 2^n$), so $2^{n+1} =
\Theta(2^n)$.

\item $T(1024) = T(2^{10}) = T(1) + 10 = 1 + 10 = 11$ (using $T(n) = 1 + \log_2 n$ from
A11).

\item \textbf{(c) $\Theta(n^2)$.} The inner loop runs $n-i$ times for each $i$, so the
total is $\sum_{i=0}^{n-1}(n-i) = \sum_{j=1}^{n} j = \frac{n(n+1)}{2} = \Theta(n^2)$ --
same order as a full nested loop, just with a smaller constant ($\approx n^2/2$ instead of
$n^2$).

\item \textbf{True.} $50$ is a constant, so $O(n^{50})$ is of the form $O(n^k)$ for fixed
$k$, which is exactly the definition of polynomial time (even though it is impractical for
real use, as discussed in A15's caveats).

\item $\lceil \log_2 17 \rceil = 5$ (since $2^4=16 < 17 \le 32 = 2^5$).

\item \textbf{True.} From A19, the total copies after $n$ appends is $2^k - 1$ where $2^k$
is the smallest power of two $\ge n$. For $n=17$, $2^k = 32$, so total copies $= 31 < 34 =
2 \times 17$.

\item $\Theta(n^2)$. \texttt{build\_women\_rank} fills an $n \times n$ table
\texttt{women\_rank[w][m]}, doing $O(1)$ work per (woman, man) pair via one pass over each
woman's preference list of length $n$, for a total of $n \times n = \Theta(n^2)$.

\end{enumerate}

\end{document}
